\section{结论与后续计划}
\label{sec:conclusion}

本文介绍了 \sysname{}，一个从年报到深度研报的可追溯投研多智能体系统。它的出发点很朴素：模型擅长读、写与判断，不擅长、也不应该负责计算和记忆数字。围绕这一分工，我们把九条设计规范写成程序可检查的不变量，构建了证据约束的双通道抽取、单一真值表驱动的估值链、按展示精度定容差的数字护栏与约束改写的数字锁；在此之上，用隔离工作区、五个上下文独立的会话与十一道闸门组织研报生产，用对生成方不可见的四层评分标准与盲评闭环组织质检，并用重复评审、程序造缺陷与人工复核三组证据校准评审模型本身。

初步结果表明：面向真实年报版式的解析让规则通道在三家公司上覆盖了 93\% 的核验值且零错配，使双通道校验在真实数据上真正起作用；先搜集资料、再做深度专题调研的流水线显著提升了研报的分析质量得分；评审模型在可由文本局部判断的缺陷上可靠，在需要通读全文判断“是否存在”的条目上仍需改进评分标准与校准方法。

决赛前的计划按优先级排列：（1）扩充真实年报与人工埋错研报评测集，报告抽取与纠错在多行业上的指标；（2）改进劣化算子（跨全文删除并同步删除图表）并完成三项设计的消融；（3）季报的规则通道表头识别（“本报告期 / 年初至报告期末”），让交叉校验覆盖季报；（4）完善终端界面与运行回放，用于现场演示时从研报中的任一数字回溯到年报页码。
